\subsection*{一步一步建立直观}

\paragraph{1. 先读圆上的坐标，再给它们命名}

把单位圆上一点投影到横轴和纵轴，两个读数就是余弦与正弦。角度决定的是位置，不是单独某条三角形边的长度。这个表示在越过直角后仍然有效，坐标符号自然改变，所以三角函数不必被限制在锐角直角三角形中。



\paragraph{2. 用对称性检查符号和周期}

转到半圈，横坐标变为负一，纵坐标回到零。继续移动点，可以通过圆的对称性推断其他象限的符号与数值关系。因为横纵坐标来自同一根单位向量，它们的平方和始终为一；这条恒等式有直接的长度解释。



\paragraph{3. 区分单位圆与一般极坐标}

加入径向增长后，点不再位于单位圆。方向仍用角度描述，坐标却要乘以当前半径，平方和也变成半径的平方。这个对照帮助避免把 $x=\cos\theta$、$y=\sin\theta$ 无条件套到任意平面点上；一般情况还必须记录模长。
